https://arxiv.org/abs/2305.13179 – another example of fine tuning to encourage probabilistic capabilities
https://arxiv.org/abs/2402.09614 – maybe a useful dataset
**https://arxiv.org/abs/2503.17523 – from Google, similar message to ours but not grounded in Zellner. They have some different datasets that we could maybe use. They did some fine-tuning, but not totally clear how (I haven’t read yet)

Concurrent work (arXiv:2503.17523) models LLM belief updating as Bayesian inference in a multi‑turn dialogue setting and fine‑tunes via SFT on oracle answers and constructed Bayesian targets, primarily on synthetic/toy tasks; unlike our approach, it is not grounded in Zellner’s information‑theoretic framework, adopts a simpler setup, and offers preliminary insights rather than a formal information‑theoretic analysis or a broader, more generalizable training method.



Here are some references that can help with our motivation and thatwe could consider using their datasets and address similar things as we do (I know we said we were freezing the datasets, but we should still consider these since we can compare to previous literature):
- Decision-Making Behavior Evaluation Framework for LLMs under Uncertain Context
- DeLLMa: Decision Making Under Uncertainty with Large Language Models
- QUITE: Quantifying Uncertainty in Natural Language Text in Bayesian Reasoning Scenarios